\section{Notation and conventions}
%\label{section2}

Throughout the paper we f\/ix an integer $n\geq 2$.
The ground f\/ield will be ${\mathbb C}$.
For $a \in {\mathbb Z}$, we write $\mathbb Z_{\geq a}$ for the set of all integers~$m$ such that $m \geq a$.
Similarly, we def\/ine $\mathbb Z_{< a}$, etc.
By~$\mathfrak{gl}(n)$ we denote the general linear Lie algebra consisting
of all $n\times n$ complex matrices, and by $\{E_{i,j}\,|\, 1\leq i,j \leq n\}$,
the standard basis of $\mathfrak{gl}(n)$ of elementary matrices.
We f\/ix the standard Cartan subalgebra $\mathfrak h$, the standard triangular decomposition and the corresponding basis
of simple roots of $\mathfrak{gl}(n)$.
The weights of $\mathfrak{gl}(n)$ will be written as~$n$-tuples $(\lambda_1,\dots,\lambda_n)$.

For a~Lie algebra ${\mathfrak a}$ by $U(\mathfrak a)$ we denote the universal enveloping algebra of ${\mathfrak a}$.
Throughout the paper $U = U(\mathfrak{gl} (n))$.
For a~commutative ring~$R$, by $\Specm R$ we denote the set of maximal ideals of~$R$.

We will write the vectors in $\mathbb{C}^{\frac{n(n+1)}{2}}$ in the following form:
\begin{gather*}
L=(l_{ij})=(l_{n1},\dots,l_{nn}| \dots|l_{21}, l_{22}|l_{11}).
\end{gather*}
For $1\leq j \leq i \leq n$, $\delta^{ij} \in {\mathbb Z}^{\frac{n(n+1)}{2}}$ is def\/ined by $(\delta^{ij})_{ij}=1$ and
all other $(\delta^{ij})_{k\ell}$ are zero.

For $i>0$ by $S_i$ we denote the~$i$th symmetric group.
Throughout the paper we set $G:=S_n\times\dots \times S_1$.

\section{Gelfand--Tsetlin modules}
\label{section3}

Recall that $U=U(\mathfrak{gl} (n))$.
Let for $m\leqslant n$, $\mathfrak{gl}_{m}$ be the Lie subalgebra of $\mathfrak{gl} (n)$ spanned by $\{E_{ij} \,|\,
i,j=1,\ldots,m \}$.
We have the following chain
\begin{gather*}
\mathfrak{gl}_1\subset \mathfrak{gl}_2\subset \cdots \subset \mathfrak{gl}_n.
\end{gather*}
It induces the chain $U_1\subset$ $U_2\subset \dots \subset U_n$ for the universal enveloping algebras
$U_{m}=U(\mathfrak{gl}_{m})$, $1\leq m\leq n$.
Let $Z_{m}$ be the center of $U_{m}$.
The subalgebra of~$U$ generated by $\{Z_m \,|\, m=1,\ldots, n \}$ will be called the \emph{$($standard$)$ Gelfand--Tsetlin
subalgebra} of~$U$ and will be denoted by ${{\Gamma}}$~\cite{DFO3}.

\begin{Definition}
%\label{definition-of-GZ-modules}
A~f\/initely generated~$U$-module~$M$ is called a~\emph{Gelfand--Tsetlin module $($with respect to ${\Gamma})$} if
\begin{gather*}
M=\bigoplus_{\mathsf m\in\Specm{\Gamma}}M(\mathsf m),
\end{gather*}
where $M(\mathsf m)=\{v\in M\,|\, \mathsf m^{k}v=0~\text{for some}~k\geq 0\}$.
\end{Definition}

For each $\mathbf m\in\Specm{\Gamma}$ we have associated a~character $\chi_{\mathbf m}:{\Gamma}\rightarrow{\Gamma}/\mathbf
m\sim\mathbb{C}$.
In the same way, for each non-zero character $\chi:{\Gamma}\rightarrow\mathbb{C}$ we have that $Ker(\chi)$ is a~maximal
ideal of ${\Gamma}$.
So, we have a~natural identif\/ication between characters of ${\Gamma}$ and elements of $\Specm{\Gamma}$.
Using characters we can def\/ine Gelfand--Tsetlin modules.
A~$U$-module~$M$ is called \emph{Gelfand--Tsetlin module} (with respect to~$\Gamma$) if
\begin{gather*}
M=\bigoplus_{\chi\in\Gamma^{*}}M(\chi),
\end{gather*}
where $M(\chi)=\{v\in M: \forall\, g\in\Gamma, \, \exists\, k\in\mathbb{Z}_{>0}~\text{such that}~(g-\chi(g))^{k}v=0 \}$.
The {\it Gelfand--Tsetlin support} of~$M$ is the set $\Supp_{\rm GT}(M):=\{\chi\in{\Gamma}^{*}:M(\chi)\neq 0\}$.

\begin{Lemma}
\label{lemma elements of Gamma separates tableaux}
Any submodule of a~Gelfand--Tsetlin module over $\mathfrak{gl}(n)$ is a~Gelfand--Tsetlin mo\-dule.
\end{Lemma}

\begin{proof}
The proof is standard, but for a~sake of completeness, we provide the important details.
Let~$M$ be a~Gelfand--Tsetlin $\mathfrak{gl}(n)$-module and~$N$ any submodule of~$M$.
We will prove that, if $\{\chi_{1},\dots,\chi_{k} \}$ is a~set of distinct Gelfand--Tsetlin characters in
$\Supp_{\rm GT}(M)$ such that $\sum\limits_{i=1}^{k}v_{i}\in N$ with $v_{i}\in M(\chi_{i})$, then $v_{i}\in N$ for all
$i=1,\ldots,k$.

Without loss of generality we assume that $k=2$.
Since $\chi_{1}\neq\chi_{2}$, there exist $g\in\Gamma$ and $r\leq s$ in $\mathbb{Z}_{\geq 0}$ such that
$\chi_{1}(g)\neq\chi_{2}(g)$, $(g-\chi_{1}(g))^{r}(v_{1})=0$ and $(g-\chi_{2}(g))^{s}(v_{2})=0$.
Let $a:=\chi_{1}(g)$ and $b:= \chi_{2}(g)$, Then, if $w=v_{1}+v_{2}$ we have $(g-b)^{s}w=(g-b)^{s}v_{1}\in N$.
Let $y:= (g-b)^sv_1$.
We have that $y \in N$ on one hand and
\begin{gather*}
y=((g-a)+(a-b))^{s}v_{1}=\sum\limits_{k=0}^{r-1} {s\choose k}(a-b)^{s-k}(g-a)^{k}v_{1}\in N
\end{gather*}
on the other.
As ${s\choose k}(a-b)^{s-k}\neq 0$ for any~$k$, using that $(g-a)^{r-1}y\in N$, we obtain \mbox{$(g-a)^{r-1}v_1\in N$}.
Reasoning in the same way, from $(g-a)^{r-i}y\in N$, and $(g-a)^{r-1}v_{1},\ldots$, $(g-a)^{r-i+1}v_{1}\in N $ we obtain $
x^{r-i}v_{1}\in N$.
Hence $v_{1}\in N$ and consequently, $v_{2}\in N$.
\end{proof}

One can choose the following generators of~$\Gamma$: $\{c_{mk} \,|\, 1\leq k\leq m\leq n \}$, where
\begin{gather}
\label{equ_3}
c_{mk}= {\sum\limits_{(i_1,\ldots,i_k)\in \{1,\ldots,m \}^k}} E_{i_1 i_2}E_{i_2 i_3}\cdots E_{i_k i_1}.
\end{gather}

Let~$\Lambda$ be the polynomial algebra in the variables $\{\lambda_{ij}\, | \, 1\leqslant j\leqslant i\leqslant n \}$.
The action of the symmetric group $S_i$ on $\{\lambda_{ij} |$ $1\leqslant j\leqslant i\}$ induces the action of $G = S_n
\times\dots \times S_1$ on~$\Lambda$.
There is a~natural embedding $\imath:{{\Gamma}}{\longrightarrow}$~$\Lambda$ given~by
$\imath(c_{mk})=\gamma_{mk}(\lambda)$, where
\begin{gather}
\label{def-gamma}
\gamma_{mk}(\lambda)=\sum\limits_{i=1}^m (\lambda_{mi}+m-1)^k \prod\limits_{j\ne i} \left(1-\frac{1}{\lambda_{mi}-\lambda_{mj}} \right).
\end{gather}
Hence,~$\Gamma$ can be identif\/ied with $G$-invariant polynomials in~$\Lambda$.

\begin{Remark}
\label{maximal ideals differ by permutations}
In what follows, we will identify the set $\Specm \Lambda$ of maximal ideals of~$\Lambda$ with the set
$\mathbb{C}^{\frac{n(n+1)}{2}}$.
Then we have a~surjective map $\pi: \Specm\Lambda \rightarrow \Specm{\Gamma}$.
Moreover, since~$\Lambda$ is integral over~$\Gamma$, there are f\/initely many maximal ideals of~$\Lambda$ that map to
a~f\/ixed maximal ideal of~$\Gamma$.
The dif\/ferent maximal ideals of~$\Lambda$ are obtained from each other under permutations in the group~$G$.
\end{Remark}
If $\pi(\ell)=\mathsf m$ for some $\ell\in \Specm \Lambda$, then we write $\ell=\ell_{\mathsf m}$ and say that
$\ell_{\mathsf m}$ is {\it lying over} $\mathsf m$.

\section[Finite-dimensional modules of $\mathfrak{gl}(n)$]{Finite-dimensional modules of $\boldsymbol{\mathfrak{gl}(n)}$}
\label{section4}

In this section we recall a~classical result of Gelfand and Tsetlin which provides an explicit basis for every
irreducible f\/inite-dimensional $\mathfrak{gl}(n)$-module.

\begin{Definition}
For a~vector $L=(l_{ij})$ in $\mathbb{C}^{\frac{n(n+1)}{2}}$, by $T(L)$ we will denote the following array with entries
$\{l_{ij}:1\leq j\leq i\leq n\}$

\begin{center}
\Stone{\mbox{\scriptsize {$l_{n1}$}}}\Stone{\mbox{\scriptsize {$l_{n2}$}}}\hspace{1cm} $\dots$ \hspace{1cm}
\Stone{\mbox{\scriptsize {$l_{n,n-1}$}}}\Stone{\mbox{\scriptsize {$l_{nn}$}}}
\\[0.2pt] \Stone{\mbox{\scriptsize {$l_{n-1,1}$}}}\hspace{1.5cm} $\dots$ \hspace{1.5cm} \Stone{\mbox{\tiny
{$l_{n-1,n-1}$}}}
\\[0.3cm] \hspace{0.2cm}$\dots$ \hspace{0.8cm} $\dots$ \hspace{0.8cm} $\dots$
\\[0.3cm] \Stone{\mbox{\scriptsize {$l_{21}$}}}\Stone{\mbox{\scriptsize {$l_{22}$}}}
\\[0.2pt] \Stone{\mbox{\scriptsize {$l_{11}$}}}
\end{center}
Such an array will be called a~\emph{Gelfand--Tsetlin tableau} of height~$n$.
A~Gelfand--Tsetlin tableau of height~$n$ is called \emph{standard} if $l_{ki}-l_{k-1,i}\in\mathbb{Z}_{\geq 0}$ and
$l_{k-1,i}-l_{k,i+1}\in\mathbb{Z}_{> 0}$ for all $1\leq i\leq k\leq n-1$.
\end{Definition}
Note that, for sake of convenience, the second condition above is slightly dif\/ferent from the original condition
in~\cite{GT}.

\begin{Theorem}[\cite{GT}]\label{Gelfand--Tsetlin theorem}
Let $L(\lambda)$ be the finite-dimensional irreducible module over $\mathfrak{gl}(n)$ of highest weight
$\lambda=(\lambda_{1},\ldots,\lambda_{n})$.
Then there exist a~basis of $L(\lambda)$ consisting of all standard tableaux $T(L) = T(l_{ij})$ with fixed top row
$l_{nj}=\lambda_j-j+1$.
Moreover, the action of the generators of~$\mathfrak{gl}(n)$ on~$L(\lambda)$ is given by the Gelfand--Tsetlin
formulas:
\begin{gather}
E_{k,k+1}(T(L)) =-\sum\limits_{i=1}^{k}\left(\frac{\prod\limits_{j=1}^{k+1}(l_{ki}-l_{k+1,j})}{\prod\limits_{j\neq
i}^{k}(l_{ki}-l_{kj})}\right)T\big(L+\delta^{ki}\big),
\nonumber
\\
E_{k+1,k}(T(L)) =\sum\limits_{i=1}^{k}\left(\frac{\prod\limits_{j=1}^{k-1}(l_{ki}-l_{k-1,j})}{\prod\limits_{j\neq
i}^{k}(l_{ki}-l_{kj})}\right)T\big(L-\delta^{ki}\big),
\nonumber
\\
E_{kk}(T(L)) =\left(k-1+\sum\limits_{i=1}^{k}l_{ki}-\sum\limits_{i=1}^{k-1}l_{k-1,i}\right)T(L),
\label{Gelfand--Tsetlin formulas}
\end{gather}
if the new tableau $T(L\pm\delta^{ki})$ is not standard, then the corresponding summand of $E_{k,k+1}(T(L))$ or
$E_{k+1,k}(T(L))$ is zero by definition.
Furthermore, for $s\leq r$,
\begin{gather}
\label{action of Gamma in finite dimensional modules}
c_{rs}(T(L))=\gamma_{rs}(l)T(L),
\end{gather}
where $\{c_{rs}\}$ are the generators of~$\Gamma$ defined in~\eqref{equ_3} and $\gamma_{rs}$ are defined
in~\eqref{def-gamma} $($see~{\rm \cite{Zh})}.
\end{Theorem}
\noindent The formulas above are called \emph{Gelfand--Tsetlin formulas} for $\mathfrak{gl}(n)$.
These formulas were extended to the case of $U_{q}(\mathfrak{gl}(n))$ in~\cite{Tar-Naz}.

\section[Generic Gelfand--Tsetlin modules of $\mathfrak{gl}(n)$]{Generic Gelfand--Tsetlin modules of $\boldsymbol{\mathfrak{gl}(n)}$}
\label{section5}

Theorem~\ref{Gelfand--Tsetlin theorem} gives an explicit realization of any irreducible f\/inite-dimensional
$\mathfrak{gl}(n)$-module.
Using the Gelfand--Tsetlin formulas, Drozd, Futorny and Ovsienko def\/ined the class of inf\/inite-dimensional generic
modules for $\mathfrak{gl}(n)$ in~\cite{DFO3}.

\begin{Definition}
\label{def-gen}
A~Gelfand--Tsetlin tableau $T(L)$ \big(equivalently, $L\in{\mathbb C}^{\frac{n(n+1)}{2}}$\big) is called \emph{generic} if
$l_{ki}-l_{kj}\notin\mathbb{Z}$ for all $1\leq i\neq j \leq k \leq n-1$.
A~character~$\chi$ and $\mathsf n = \Ker \chi$ are called \emph{generic} if $\ell_{\mathsf n}$ is generic for one choice
(hence for all choices) of $\ell_{\mathsf n}$ lying over $\mathsf n$.
A~Gelfand--Tsetlin module~$M$ will be called a~\emph{generic Gelfand--Tsetlin module} if every $\mathsf n$ in
$\Supp_{\rm GT}(M)$ is generic.
\end{Definition}

\begin{Theorem}[\protect{\cite[Section 2.3]{DFO3} and~\cite[Theorem 2]{Maz2}}]%\label{Generic Gelfand--Tsetlin modules}
Let $T(L) = T(l_{ij})$ be a~generic Gelfand--Tsetlin tableau of height~$n$.
Denote by ${\mathcal B}(T(L))$ the set of all Gelfand--Tsetlin tableaux $T(R) = T(r_{ij})$ satisfying $r_{nj}=l_{nj}$,
$r_{ij}-l_{ij}\in\mathbb{Z}$ for
$1\leq j\leq i \leq n-1$.
\begin{itemize}\itemsep=0pt
\item[$(i)$] The vector space $V(T(L)) = \Span {\mathcal B}(T(L))$ has a~structure of a~$\mathfrak{gl}(n)$-module with
action of the generators of $\mathfrak{gl}(n)$ given by the Gelfand--Tsetlin formulas~\eqref{Gelfand--Tsetlin formulas}.
\item[$(ii)$] The action of the generators of~$\Gamma$ on the basis elements of $V(T(L))$ is given by~\eqref{action of Gamma in finite dimensional modules}.
\item[$(iii)$] The $\mathfrak{gl}(n)$-module $V(T(L))$ is a~Gelfand--Tsetlin module all of whose Gelfand--Tsetlin
multiplicities are~$1$.
\end{itemize}
\end{Theorem}

\begin{Remark}

The basis of the module in the previous theorem is
\begin{gather*}
{\mathcal B}(T(L)) = \big\{T(L+z):z\in\mathbb{Z}^{\frac{n(n+1)}{2}}~\text{and}~z_{n1}=\cdots=z_{nn}=0\big\}.
\end{gather*}
By a~slight abuse of notation we will identify elements in $\mathbb{Z}^{\frac{n(n-1)}{2}}$ with elements
$z\in\mathbb{Z}^{\frac{n(n+1)}{2}}$ such that $z_{n1}=\cdots=z_{nn}=0$.
This will allow us to write $T(L+z)$ for $z\in\mathbb{Z}^{\frac{n(n-1)}{2}}$.
\end{Remark}

\begin{Remark}
\label{elements of Gamma separates tableaux}
In what follows, we will apply Lemma~\ref{lemma elements of Gamma separates tableaux} and use that the elements
of~$\Gamma$ separate the tableaux in the submodules of $V(T(L))$ in the following sense.
Let~$N$ be a~$\mathfrak{gl}(n)$-submodule of $V(T(L))$, $g \in \mathfrak{gl}(n)$, and $T(R)$ be a~tableau in~$N$.
Then, if $g \cdot T(R) = \sum\limits_i c_i T(R_i)$ for some distinct tableaux $T(R_i)$ in ${\mathcal B}(T(L))$ and
nonzero $c_i \in {\mathbb C}$, we have $T(R_i) \in N$ for all~$i$.
\end{Remark}

\begin{Theorem}
\label{uniqueness in generic case}
If $\mathsf n\in\Specm\Gamma$ is generic, then there exists a~unique irreducible Gelfand--Tsetlin module~$N$ such that
$N(\mathsf n)\neq 0$.
\end{Theorem}
\begin{proof}
Let $X_{\mathsf n}=U/U \mathsf n$.
We know that $X_{\mathsf n}=U/U \mathsf n$ is a~Gelfand--Tsetlin module.
Furthermore, any irreducible Gelfand--Tsetlin module~$M$ with $M(\mathsf n)\neq 0$ is a~homomorphic image of $X_{\mathsf
n}$, and~$X_{\mathsf n}(\mathsf n)$ maps onto~$M(\mathsf n)$.
Since both spaces $X_{\mathsf n}(\mathsf n)$ and $M(\mathsf n)$ are~$\Gamma$-modules then the projection $X_{\mathsf
n}(\mathsf n) \to M(\mathsf n)$ is a~homomorphism of~$\Gamma$-modules (see also~\cite[Corollary~5.3]{FO2}).
Taking into account that $\dim X_{\mathsf n}(\mathsf n) \leq 1$, we conclude that $X_{\mathsf n}$ has a~unique maximal
submodule (which does not intersect $X_{\mathsf n}(\mathsf n)$) and hence there exist a~unique irreducible module~$N$
with $N(\mathsf n)\neq 0$.
\end{proof}

\begin{Definition}
\label{def-cont}
If $T(R)$ is a~generic tableau and $\mathsf r\in\Specm\Gamma$ corresponds to~$R$ then, the unique module~$N$ such that
$N(\mathsf r) \neq 0$ is called the \emph{irreducible Gelfand--Tsetlin module containing $T(R)$}, or simply, the
\emph{irreducible module containing $T(R)$}.
\end{Definition}

Our goal is to describe explicitly the irreducible Gelfand--Tsetlin module containing $T(R)$ for every generic tableau $T(R)$.
